\section{Introduction}
\label{sec:intro}

Progress in end-to-end autonomous driving is measured by benchmarks.
Real-world closed-loop testing is unavailable at research scale, so policies are ranked by proxy
evaluators: displacement from a logged human trajectory~\cite{zhai2023admlp,li2024egostatus},
non-reactive pseudo-simulation such as NAVSIM~\cite{dauner2024navsim}, two-stage
pseudo-simulation~\cite{cao2025pseudosim}, and closed-loop simulators such as
CARLA-based Bench2Drive~\cite{jia2024bench2drive} or BridgeSim~\cite{zhao2026bridgesim}.
Papers routinely report gains of one or two points on these scores.
Whether such a gain means better driving depends on a property of the evaluator that is rarely
measured: does it rank policies in their true order?

Existing studies answer a weaker question.
They correlate the published numbers of a handful of methods across benchmarks~\cite{wang2026olcl},
or compare one proxy with one simulator~\cite{cao2025pseudosim,zhao2026bridgesim}.
In both cases evaluator is compared with evaluator; no study has a ground truth, so a disagreement
cannot say which side is wrong, and an agreement can be shared bias.

We remove this obstacle by constructing policies whose order is known.
A pretrained policy degraded in graded steps---its plan delayed, drifted sideways, its cameras
blanked---must drive worse as the step grows.
An evaluator that does not recover this order is blind to that failure.
Two safeguards make the construction usable.
First, the assumption ``more degradation is worse'' is itself tested: a degradation family is
admitted only if a closed-loop reference ranks its levels with statistical confidence, and,
to avoid circularity, this selection is made on one half of the scenarios while every evaluator,
the reference included, is measured on the other half.
Second, sensitivity to harmful change is not enough: an evaluator that penalises every change to the
plan passes an ordering test while mis-ranking real policies.
We therefore add a \emph{specificity arm} of benign perturbations---the same edits applied only to
the part of a plan that a replanning controller never executes---which a correct evaluator must
leave unpenalised.

We apply this protocol to five evaluators that can be run on the same scenarios: open-loop average
displacement error, NAVSIM v1 PDMS, NAVSIM v2 EPDMS in one stage and in its two-stage
pseudo-simulation, and BridgeSim, a closed-loop simulator that runs the same checkpoints and serves
as the reference.
The policy population comprises TransFuser~\cite{chitta2023transfuser}, Latent TransFuser and the
Ego Status MLP of NAVSIM, three training seeds each, with six degradation families at three levels
and two benign families.

\paragraph{Contributions.}
\begin{itemize}
\item A ground-truth ordering test for driving evaluators, with the ordering itself validated on
held-out scenarios before use (\cref{sec:method}).
\item A specificity test that separates evaluators which detect bad driving from those which
detect any change to the plan (\cref{sec:method:spec}).
\item A controlled study of five evaluators on shared scenarios, with paired bootstrap intervals,
seed replication and a closed-loop test-retest ceiling (\cref{sec:results}).
\item An open harness and all per-scenario scores.
\end{itemize}
